\begin{proposition}[Finite limits and the cusp limit]\hypertarget{tag-01RF}{}\label{fam:src:F.3.2}\label{fam:st:F.3.2}
\StatementRecord{fam:src:F.3.2}{proved}{01RF}{\SeriesRef{fam:src:F.3.1}}

At the two finite punctures the period map has the holomorphic
finite-cover limits
\[
 \Pi_1=
 \begin{pmatrix}2(2-\rho)&\rho&1&0\\
 \beta_c(z_1)&(2-\rho)/3&0&1\end{pmatrix},\qquad
 \Pi_2=
 \begin{pmatrix}3(1-i)&i&1&0\\
 \beta_c(z_2)&(1-i)/2&0&1\end{pmatrix},
 \quad \rho=e^{\pi i/3}.
\]
Their monodromies are the given order-three and order-four matrices.
At the cusp the curve leaves every compact subset of
$\Gamma\backslash\mathcal O_+$.  On the prescribed cusp branch its
Hodge plane in the compact dual tends to
$\mathbb C\gamma+\mathbb Cu$.  The latter plane has complex dimension
two in common with its conjugate and therefore does not represent a
point of $\mathcal D$.
\end{proposition}



\begin{proof}
\begin{enumerate}[label=\arabic*.,leftmargin=*]
\item The finite values of $\tau,\mu$ follow by evaluating their transformation
laws at the fixed points: $\tau(z_1)=\rho$,
$\mu(z_1)=(2-\rho)/3$, $\tau(z_2)=i$,
$\mu(z_2)=(1-i)/2$.  Substitution in $\Pi$ gives the two matrices.
The finite covering comparison and its descent were proved in \ref{fam:src:F.3.1}.
The unipotent cusp transformation is
$\tau\mapsto\tau-1$, $\mu\mapsto\mu$, $\beta\mapsto\beta+1$.
On the cusp branch, $\operatorname{Im}\tau\to+\infty$,
$\mu\to\mu_0$, and $\beta_c+\tau\to b_0+c$, from the holomorphic
cusp expansions.  Consequently
\[
 D_c=-\operatorname{Im}\tau+O(1),\qquad
 q_e(\Pi_c)\longrightarrow0^+,
 \qquad j_\Gamma(p_c(t))=1728t\longrightarrow\infty.
\]
The last equality alone excludes a convergent subnet in any compact
subset of the proper quotient; the first two identities give the
position relative to the maximal proper-action open set.
The rows of $\Pi_c$, regarded as covectors, are
$6\mu\gamma+\tau u+w$ and $\beta_c\gamma+\mu u+\delta$.
Dividing each by $\tau$ gives limits $u$ and $-\gamma$.
Their exterior product has a nonzero limit, so their span converges
in the Grassmannian to the stated plane.  It is spanned by real
covectors, hence equals its conjugate.  A complex structure has
$H^{1,0}\cap H^{0,1}=0$; here the intersection is the entire
two-dimensional plane.
\end{enumerate}
\end{proof}